\chapter{AWS platform}\label{ch:aws}
\chapterstatus{complete}{2026-11-16}{\src{infra/aws/org}, \src{infra/aws/app}, \src{Dockerfile}}

Slotty runs in its own AWS account, \code{slotty-prod}, in us-east-1. Everything in this chapter
is created by Terraform from \src{infra/aws}; nothing was built by hand except the two steps AWS
requires a person for (the SES production-access request and the DNS delegation from
Cloudflare). \Cref{fig:aws} shows the whole deployment.

\begin{figure}[p]
  \centering
  \makebox[\linewidth][c]{% AWS deployment (chapter 8), mirroring infra/aws/app. Grid: columns 25 mm apart, rows in mm.
% Solid: requests and data. Dashed: egress and events. Dotted: deployment.
% Not drawn (to keep the picture readable; see the text): image pulls through the S3 endpoint,
% secrets from SSM at task start, logs to CloudWatch.
\begin{tikzpicture}[x=1mm, y=1mm, every node/.append style={font=\sffamily\scriptsize},
    n/.style={box=20mm, minimum height=9mm, inner sep=2pt},
    sub/.style={font=\sffamily\tiny, text=muted}]
  % ── top row: users and the edge ──
  \node[actor, text width=17mm] (users) at (0,31) {Users};
  \node[aws=network, n] (r53) at (25,0) {Route 53\\[-2pt]\textcolor{muted}{\tiny DNS zone}};
  \node[aws=network, n] (cf) at (50,0) {CloudFront\\[-2pt]\textcolor{muted}{\tiny TLS, cache}};
  \node[aws=network, n] (api) at (75,0) {API Gateway\\[-2pt]\textcolor{muted}{\tiny HTTP API}};
  \node[aws=security, n] (acm) at (50,14) {ACM\\[-2pt]\textcolor{muted}{\tiny certificate}};
  \node[ext, text width=20mm] (gh) at (125,31) {GitHub Actions\\[-2pt]\textcolor{muted}{\tiny OIDC}};
  % ── VPC ──
  \node[aws=compute, n] (nat) at (0,-28) {fck-nat\\[-2pt]\textcolor{muted}{\tiny t4g.nano}};
  \node[aws=network, n] (s3e) at (25,-28) {S3 endpoint\\[-2pt]\textcolor{muted}{\tiny gateway}};
  \node[aws=network, n] (vpcl) at (75,-44) {VPC link};
  \node[aws=compute, n] (weba) at (0,-61) {web\\[-2pt]\textcolor{muted}{\tiny Fargate, AZ a}};
  \node[aws=compute, n] (webb) at (25,-61) {web\\[-2pt]\textcolor{muted}{\tiny Spot, AZ b}};
  \node[aws=compute, n] (wrk) at (50,-61) {worker\\[-2pt]\textcolor{muted}{\tiny Spot}};
  \node[aws=compute, n] (mig) at (75,-61) {migrate\\[-2pt]\textcolor{muted}{\tiny one-off}};
  \node[aws=database, box=34mm, minimum height=9mm, inner sep=2pt] (rds) at (37.5,-83) {RDS PostgreSQL 17\\[-2pt]\textcolor{muted}{\tiny encrypted, TLS only}};
  % ── regional services ──
  \node[aws=network, n] (cm) at (100,-44) {Cloud Map\\[-2pt]\textcolor{muted}{\tiny healthy tasks}};
  \node[aws=compute, n] (ecr) at (125,-18) {ECR\\[-2pt]\textcolor{muted}{\tiny images}};
  \node[aws=security, n] (ssm) at (125,-36) {SSM\\[-2pt]\textcolor{muted}{\tiny secrets}};
  \node[aws=manage, n] (cw) at (125,-63) {CloudWatch\\[-2pt]\textcolor{muted}{\tiny logs, alarms}};
  \node[aws=business, n] (ses) at (100,-83) {SES\\[-2pt]\textcolor{muted}{\tiny email}};
  \node[aws=integration, n] (sns) at (125,-83) {SNS\\[-2pt]\textcolor{muted}{\tiny alerts}};

  \begin{scope}[on background layer]
    \node[subnet=public, fit=(nat), inner xsep=3mm, inner ysep=4mm] (pub) {};
    \node[subnet=private, fit=(vpcl)(weba)(mig), inner xsep=3mm, inner ysep=4.5mm] (app) {};
    \node[subnet=data, fit=(rds), inner xsep=3mm, inner ysep=4mm] (data) {};
    \node[vpc, fit=(pub)(s3e)(app)(data), inner sep=2.5mm] (vpc) {};
    \node[account, fit=(r53)(acm)(vpc)(cm)(ssm)(sns), inner sep=3.5mm] (acct) {};
  \end{scope}
  \node[grouplabel, text=awsgeneric] at (acct.north west) {Account slotty-prod, us-east-1};
  \node[grouplabel, text=awsnetwork] at (vpc.north west) {VPC 10.0.0.0/16};
  \node[sub, anchor=north west] at (pub.north west) {public};
  \node[sub, anchor=north west] at (app.north west) {app subnets (2 AZs)};
  \node[sub, anchor=north west] at (data.north west) {db subnets};

  % request path
  \draw[flow] (users.south) |- (r53.west);
  \draw[flow] (users.east) -| ([xshift=-6mm]cf.north);
  \draw[flow] (cf) -- (api);
  \draw[ctrl] (acm) -- (cf);
  \draw[flow] (api) -- (vpcl);
  \draw[flow] (vpcl) -- (cm);
  \draw[flow] (vpcl.west) -| (weba.north);
  \draw[flow] (vpcl.west) -| (webb.north);
  % data
  \draw[flow] (weba.south) |- (rds.west);
  \draw[flow] (webb.south) -- (webb.south |- rds.north);
  \draw[flow] (wrk.south) -- (wrk.south |- rds.north);
  \draw[flow] (mig.south) |- (rds.east);
  % egress: the app subnets' default route goes through the NAT
  \draw[async] (nat.south |- app.north) -- node[sub, right]{default route} (nat.south);
  \draw[async] (app.east |- wrk) -| node[sub, pos=.25, above]{SMTP, via NAT} (ses.north);
  \draw[async] (ses) -- node[sub, above]{bounces} (sns);
  \draw[async] (cw) -- (sns);
  % deployment
  \draw[ctrl] (gh) -- node[sub, right]{push} (ecr);
  \draw[ctrl] (gh.west) -| node[sub, pos=.88, right, align=left]{run task,\\update service} ([xshift=-3mm]app.north east);
\end{tikzpicture}
}
  \caption{The AWS deployment. Solid arrows: requests and data. Dashed: egress and events.
    Dotted: deployment and configuration.}
  \label{fig:aws}
\end{figure}

\section{Accounts and guardrails}\label{sec:aws-guardrails}
An AWS Organization separates Slotty from the other workloads of the management account.
\src{infra/aws/org} manages it, applied with the management account's credentials:

\begin{xltabular}{\linewidth}{@{}P{34mm} L@{}}
  \toprule
  \thead{Resource} & \thead{What it does}\\
  \midrule
  Organization & All features, service control policies (SCPs) enabled. Protected with \code{prevent_destroy}: it can only be deleted after its member accounts are closed.\\
  \code{Workloads} OU & Holds \code{slotty-prod}; the SCPs attach here.\\
  \code{slotty-prod} account & Created by Terraform; administered by assuming \code{OrganizationAccountAccessRole} from the management account. Its ID is the \code{slotty_account_id} output.\\
  SCP \code{slotty-region-lock} & Denies every action outside us-east-1, except global services (IAM, Organizations, Route 53, CloudFront, ACM, WAF, billing, support).\\
  SCP \code{slotty-cost-guardrails} & EC2 only \code{t4g.nano/micro/small}; RDS only \code{db.t4g.micro/small}; denies Aurora, EKS, SageMaker, Redshift, ElastiCache, OpenSearch, MSK, Bedrock, NAT gateways, and every long-term purchase (Savings Plans, reservations, domain registration).\\
  SCP \code{slotty-protect} & Denies leaving the Organization, closing the account, any root-user action, and stopping or changing CloudTrail.\\
  Budget \code{slotty-prod-monthly} & \$50 a month, filtered to the account. Emails at 90\,\% and 100\,\% of actual spend and 100\,\% of forecast.\\
  Cost anomaly monitor & Daily email when an anomaly's impact reaches \$5.\\
  Organization CloudTrail & Management events from every account, to an encrypted, private S3 bucket in the management account, kept 90 days, with log-file validation.\\
  State bucket & \code{slotty-tfstate-<mgmt account>}: versioned, encrypted, TLS-only, public access blocked; native S3 locking. Holds both stacks' state.\\
  \bottomrule
\end{xltabular}

\begin{auditnote}
  Confirm the guardrails from the management account:
  \code{aws organizations list-policies-for-target --target-id <Workloads OU> --filter SERVICE_CONTROL_POLICY}
  lists the three SCPs; \code{aws budgets describe-budgets} shows the budget. From
  \code{slotty-prod}, try \code{aws ec2 run-instances --instance-type m5.large …} or any call in
  another region: both are denied by an SCP.
\end{auditnote}

\begin{risk}
  The SCPs also block things Slotty will want when it grows (a NAT gateway, a larger database).
  Raising a limit is a reviewed change to \src{infra/aws/org/scp.tf}, on purpose.
\end{risk}

\section{Request path}\label{sec:aws-request-path}
\begin{enumerate}
  \item \textbf{Route 53} answers for \code{slotty.msaeedsaeedi.com}, a subdomain delegated from
    the parent zone at Cloudflare by four NS records. A CAA record allows only Amazon to issue
    certificates for it.
  \item \textbf{CloudFront} terminates TLS (ACM certificate, TLS 1.2 minimum, HTTP/2 and HTTP/3),
    redirects HTTP to HTTPS and adds HSTS (two years), \code{nosniff}, \code{X-Frame-Options: DENY}
    and a strict referrer policy. \route{/_next/static/*} is cached at the edge (content-hashed,
    immutable); everything else passes through uncached with all viewer headers and cookies except
    \code{Host}. CloudFront adds the secret header \code{x-origin-verify} to every origin request.
  \item \textbf{API Gateway} (HTTP API) is the load balancer. Stage throttling (50 requests per
    second, bursts of 100) is the edge rate limit. Its \code{$default} route forwards everything
    through a \textbf{VPC link} into the app subnets.
  \item \textbf{Cloud Map} lists the healthy web tasks (SRV records registered by ECS); API Gateway
    picks one.
  \item The \textbf{web task} checks the secret header in \src{src/proxy.ts}, restores the public
    host for Next's Server Actions check, and serves the request.
\end{enumerate}

\begin{decision}{ADR-002}
  API Gateway instead of an Application Load Balancer: about \$1 a month at bootstrap volume
  instead of about \$18. Its limits (30\,s per request, 10\,MB responses, buffered rather than
  streamed responses) do not affect Slotty today.
\end{decision}

\begin{auditnote}
  \code{curl -I https://slotty.msaeedsaeedi.com/login} returns 200 with
  \code{strict-transport-security}; the API Gateway URL (output \code{api_endpoint}) returns 403
  for \route{/login}, because it lacks the secret header.
\end{auditnote}

\section{Network}
\begin{xltabular}{\linewidth}{@{}P{30mm} P{34mm} L@{}}
  \toprule
  \thead{Subnets (×2 AZs)} & \thead{CIDR} & \thead{Contents and routes}\\
  \midrule
  Public & \code{10.0.0.0/24}, \code{10.0.1.0/24} & The NAT instance only. Default route to the internet gateway.\\
  App (private) & \code{10.0.10.0/24}, \code{10.0.11.0/24} & ECS tasks and the VPC link. Default route through the NAT; S3 gateway endpoint for image layers.\\
  DB (private) & \code{10.0.20.0/24}, \code{10.0.21.0/24} & RDS. No default route at all.\\
  \bottomrule
\end{xltabular}

Security groups allow only the next hop: the VPC link may reach web tasks on port 3000; web and
worker tasks may reach RDS on 5432 (as may the NAT instance, for administrators' port forwarding);
the NAT instance accepts only the app subnets. The VPC's default security group has no rules.

\begin{decision}{ADR-003}
  Egress through a NAT \emph{instance} (fck-nat on a \code{t4g.nano}, about \$7 a month) instead of a
  managed NAT gateway (about \$33 plus traffic). It is a single instance with EC2 automatic
  recovery; while it is down, running tasks keep serving users but cannot send email or pull
  images. \tfvar{nat_mode = "gateway"} switches to a managed gateway (after lifting the SCP deny).
\end{decision}

\section{Compute}\label{sec:aws-compute}
ECS on Fargate, ARM64 (Graviton), with Container Insights off to save cost.

\begin{xltabular}{\linewidth}{@{}P{22mm} P{30mm} L@{}}
  \toprule
  \thead{Task} & \thead{Size} & \thead{Behaviour}\\
  \midrule
  web & 0.5 vCPU, 1 GB & Service, at least \tfvar{web_min} (2) tasks across both AZs, at most \tfvar{web_max} (4). The first task runs on regular Fargate, the rest on Fargate Spot. Scales on average CPU (target 60\,\%). Container health check on \route{/health/live}.\\
  worker & 0.25 vCPU, 0.5 GB & Service, one task on Fargate Spot. Interruptions are safe: it drains on \code{SIGTERM} within 25\,s and delivery is at least once.\\
  migrate & 0.25 vCPU, 0.5 GB & One-off task run by CI before each rollout (\cref{ch:delivery}); also used for administrative one-offs such as creating the first admin.\\
  \bottomrule
\end{xltabular}

Rollouts keep 100\,\% of web capacity (up to 200\,\% during a deploy) and use the ECS deployment
circuit breaker with automatic rollback. ECS Exec is enabled for debugging.

\section{Database}
\begin{xltabular}{\linewidth}{@{}P{38mm} L@{}}
  \toprule
  \thead{Setting} & \thead{Value}\\
  \midrule
  Engine and size & PostgreSQL 17, \code{db.t4g.micro}, single AZ (\tfvar{db_multi_az} switches)\\
  Storage & gp3, 20 GB, autoscaling to 50 GB, encrypted with the AWS-managed key\\
  Backups & Automated daily, 7-day point-in-time recovery, final snapshot on delete, deletion protection\\
  Parameters & \code{rds.force_ssl = 1}; slow statements over 1\,s logged; idle-in-transaction sessions ended after 60\,s\\
  Access & Private subnets only, not publicly accessible; credentials in SSM\\
  \bottomrule
\end{xltabular}

TLS is verified in full. The image carries the RDS CA bundle (\code{/app/certs/rds-global-bundle.pem},
checked against a pinned SHA-256 at build). The app and worker connect with
\code{sslmode=verify-full}, trusting the bundle through \envvar{NODE_EXTRA_CA_CERTS}. Prisma's
migration engine reads only the first certificate of an \code{sslcert} file, so the migrate image
instead points OpenSSL's trust store at the bundle (\envvar{SSL_CERT_FILE}) and connects with
\code{sslmode=require&sslaccept=strict}, which verifies the chain and the host name.

\begin{decision}{ADR-004}
  A single-AZ \code{db.t4g.micro} at bootstrap. A zone outage means downtime until RDS recovers
  or the database is restored elsewhere; Multi-AZ doubles the database cost. Revisit when a
  demo day's downtime would cost more than about \$14 a month.
\end{decision}

\section{Email}\label{sec:aws-email}
\aws{Amazon SES} sends from \code{no-reply@slotty.msaeedsaeedi.com} through the worker's existing
SMTP transport (port 587, STARTTLS). The domain identity has Easy DKIM (three CNAME records), a
custom MAIL FROM domain \code{mail.slotty.msaeedsaeedi.com} with SPF, and a DMARC policy of
\code{p=quarantine} with strict DKIM alignment. A configuration set, the domain's default,
publishes bounces, complaints and rejections to the alerts SNS topic; the account-level
suppression list stops further sends to bounced or complaining addresses. The SMTP credentials
belong to an IAM user that may only call \code{ses:SendRawEmail} with a \code{From} address on the
domain.

\begin{note}
  The account was created in the SES sandbox, which only delivers to verified addresses.
  Production access was requested on 2026-11-16; check \code{aws sesv2 get-account} for
  \code{ProductionAccessEnabled}.
\end{note}

\subsection{Inbound: the support address}
\code{support@slotty.msaeedsaeedi.com} is the only address that receives mail
(\src{infra/aws/app/support-mail.tf}, ADR-021). An MX record on the domain points to SES
receiving. One active receipt rule set, \code{slotty}, has one rule for that recipient, with spam
and virus scanning: it stores the raw message in a private, TLS-only S3 bucket
(\code{slotty-support-mail-<account>}, objects expire after 30 days), then invokes the
\aws{AWS Lambda} function \code{slotty-support-forwarder} asynchronously. The function drops
messages that failed the spam or virus scan, rewrites \code{From} to
\code{"<sender> via Slotty support" <support@…>} with the sender in \code{Reply-To}, and sends
it through SES to \tfvar{support_forward_to} (set in the git-ignored \code{terraform.tfvars}).
Re-sending from the domain keeps DKIM and DMARC aligned at the destination. The function may only
read that bucket prefix and send as \code{support@}. After two failed retries, the event goes to
the alerts topic. An empty \tfvar{support_forward_to} removes all of it.

\section{Secrets and configuration}
Runtime secrets are SSM Parameter Store \code{SecureString}s under \code{/slotty/}, encrypted with
the AWS-managed key and injected by ECS at task start: \envvar{DATABASE_URL},
\envvar{MIGRATE_DATABASE_URL}, \envvar{ORIGIN_VERIFY_SECRET}, \envvar{SMTP_USER},
\envvar{SMTP_PASS}, \envvar{VAPID_PUBLIC_KEY}, \envvar{VAPID_PRIVATE_KEY}, \envvar{ANALYTICS_SECRET}
(the key for the pseudonymous visitor code, web tasks only; ADR-015). Terraform generates
them, so they are also in the Terraform state; the state bucket is private and encrypted. Only the
ECS execution role may read them. Non-secret settings are plain environment variables in the task
definitions (\cref{app:config}); \src{infra/aws/deploy.sh} sets \envvar{APP_VERSION} to the
image tag on every revision it registers.

\section{Observability}
Container logs go to CloudWatch Logs (\code{/slotty/web}, \code{/slotty/worker},
\code{/slotty/migrate}, \code{/slotty/api}), kept 14 days. Ten alarms (the free allowance) email
the alerts topic:

\begin{xltabular}{\linewidth}{@{}P{42mm} L@{}}
  \toprule
  \thead{Alarm} & \thead{Fires when}\\
  \midrule
  \code{slotty-web-down}, \code{slotty-worker-down} & The service reports no CPU samples for 10 minutes (no task running)\\
  \code{slotty-web-cpu-high} & Web CPU above 85\,\% for 15 minutes despite autoscaling\\
  \code{slotty-web-memory-high} & Web memory above 85\,\% for 10 minutes\\
  \code{slotty-5xx} & More than 10 server errors in 5 minutes at API Gateway\\
  \code{slotty-db-cpu-high} & Database CPU above 80\,\% for 15 minutes\\
  \code{slotty-db-storage-low} & Less than 2 GB free\\
  \code{slotty-db-connections-high} & More than 70 connections\\
  \code{slotty-nat-status-failed} & The NAT instance fails its status checks\\
  \code{slotty-worker-failures} & A worker job crashes, or 10 or more email sends fail within 30 minutes\\
  \bottomrule
\end{xltabular}

\section{Terraform}
\begin{xltabular}{\linewidth}{@{}P{36mm} L@{}}
  \toprule
  \thead{Path} & \thead{Contents}\\
  \midrule
  \src{infra/aws/org} & Organization, account, SCPs, budget, anomaly alerts, CloudTrail, state bucket. Applied with management-account credentials.\\
  \src{infra/aws/app} & Everything inside \code{slotty-prod}, one file per concern (\code{network.tf}, \code{nat.tf}, \code{dns.tf}, \code{edge.tf}, \code{api.tf}, \code{discovery.tf}, \code{ecs.tf}, \code{rds.tf}, \code{ses.tf}, \code{support-mail.tf}, \code{secrets.tf}, \code{iam.tf}, \code{observability.tf}, \code{security.tf}). Its provider assumes \code{OrganizationAccountAccessRole}.\\
  \src{infra/aws/deploy.sh} & The rollout script CI runs (\cref{ch:delivery}).\\
  \bottomrule
\end{xltabular}

Variables and their defaults are in \src{infra/aws/app/variables.tf}; per-environment values in
\code{terraform.tfvars} (not committed). Every apply is planned first and the plan reviewed.
ECS services ignore changes to their task definition: CI owns the running image, so an
environment change made in Terraform reaches the tasks on the next deploy.

\begin{auditnote}
  \code{terraform -chdir=infra/aws/app plan -detailed-exitcode} exits 0 when the account matches
  the code. Any other exit code is drift to explain.
\end{auditnote}
